\section*{РЕФЕРАТ}

Пояснювальна записка до курсової роботи: \pageref{LastPage}~с., \total{figs}~рис.,
\total{tabs}~табл., 3~додатки, \total{refs}~джерело.

\textit{Об'єкт дослідження} -- паралельна побудова інвертованого індексу та паралельна обробка
пошукових запитів у пошуковій системі з розширенням запитів великою мовною моделлю.

\textit{Мета роботи} -- розробити контейнеризовану пошукову систему мовою Java без сторонніх
бібліотек (власний пул потоків, власна потокобезпечна структура індексу, власний HTTP/1.1- і
WebSocket-сервер, розширення запитів локальною моделлю qwen2.5:1.5b через Ollama) та
експериментально дослідити її масштабованість.

\textit{Методи дослідження}: вимірювання часу побудови індексу на 1--16 потоках для двох
реалізацій спільної структури (неблокуючої з хеш-партиціями і блокувальної з lock striping) і
для 1, 8, 64 та 256 партицій, по три прогони з випадковим чергуванням конфігурацій; аналіз
прискорення за законом Амдала та метрикою Карпа -- Флатта; розклад часу за фазами;
порівняння двох версій токенізації; навантажувальне тестування пошуку для 1--64 одночасних
клієнтів; вимірювання пам'яті під час повторних перебудов; модульне та наскрізне тестування.

\textit{Результати}. На корпусі з 16\,793 документів (1,62~ГБ) медіанний час побудови індексу
скоротився з 10,12~с на одному потоці до 2,09~с на 10 потоках (прискорення 4,85) на
процесорі з 6 ядрами і 12 апаратними потоками. Заміна токенізації, що створювала рядок на
кожне слово, байтовою хеш-таблицею задачі скоротила час у 1,28--1,48 раза за будь-якої
кількості потоків. З однією партицією блокувальний варіант на 12 потоках сумарно очікує на
блокування 22,4~с, а неблокуючий не очікує, але обидва впираються в послідовне злиття
(3,4--3,5~с); від 8 партицій синхронізація перестає бути вузьким місцем. Пропускна здатність
пошуку сягає 12\,004 запитів/с на 16 клієнтах; з розширенням через LLM вона обмежена моделлю
на рівні 2,3--2,4 запиту/с, а кешування розширень повертає її до 8\,634 запитів/с.

\bigskip
\noindent ПАРАЛЕЛЬНІ ОБЧИСЛЕННЯ, ПУЛ ПОТОКІВ, ІНВЕРТОВАНИЙ ІНДЕКС, НЕБЛОКУЮЧА СИНХРОНІЗАЦІЯ,
CAS, ЗАКОН АМДАЛА, МЕТРИКА КАРПА -- ФЛАТТА, BM25, WEBSOCKET, LLM, OLLAMA, DOCKER

\clearpage
